\documentclass[11pt,a4paper]{article}
\usepackage[utf8]{inputenc}
\usepackage[T1]{fontenc}
\usepackage[french]{babel}
\usepackage{amsmath,amssymb,amsthm}
\usepackage{amsfonts}
\usepackage{mathtools}
\usepackage{geometry}
\usepackage{hyperref}
\usepackage{enumitem}
\usepackage{booktabs}
\usepackage{array}
\usepackage{mathrsfs}
\usepackage{calrsfs}
\usepackage{xcolor}
\usepackage{microtype}

\geometry{a4paper, top=2.5cm, bottom=2.5cm, left=2.5cm, right=2.5cm}

\theoremstyle{plain}
\newtheorem{theorem}{Th\'eor\`eme}
\newtheorem{lemma}[theorem]{Lemme}
\newtheorem{proposition}[theorem]{Proposition}
\newtheorem{corollary}[theorem]{Corollaire}
\newtheorem{conjecture}[theorem]{Conjecture}

\theoremstyle{definition}
\newtheorem{definition}[theorem]{D\'efinition}
\newtheorem{example}[theorem]{Exemple}
\newtheorem{remark}[theorem]{Remarque}

\DeclareMathOperator{\Ext}{Ext}
\DeclareMathOperator{\Hom}{Hom}
\DeclareMathOperator{\Ker}{Ker}
\DeclareMathOperator{\Coker}{Coker}
% \DeclareMathOperator{\Im}{Im}  % Already defined
\DeclareMathOperator{\Spec}{Spec}
\DeclareMathOperator{\Proj}{Proj}
\DeclareMathOperator{\Gal}{Gal}
\DeclareMathOperator{\Aut}{Aut}
\DeclareMathOperator{\End}{End}
\DeclareMathOperator{\CSpin}{CSpin}
\DeclareMathOperator{\SL}{SL}
\DeclareMathOperator{\GL}{GL}
\DeclareMathOperator{\SO}{SO}
\DeclareMathOperator{\Ogrp}{O}
\DeclareMathOperator{\Spin}{Spin}
\DeclareMathOperator{\Cokername}{Coker}

\title{\textbf{Equivalence num\'erique, \'equivalence cohomologique,\\
et th\'eorie de Lefschetz des vari\'et\'es ab\'eliennes sur les corps finis}}
\author{P. Deligne (r\'edig\'e par L. Clozel)\thanks{R\'edaction bas\'ee sur une note de Deligne et une lettre du 12 juillet 1997.}}
\date{}

\begin{document}

\maketitle

\begin{center}
\textit{\`A Jean-Pierre Serre, en t\'emoignage de notre admiration}
\end{center}

\begin{abstract}\noindent
We determine the group of automorphisms generated, in the endomorphisms of the cohomology of an Abelian variety on a finite field, by the Lefschetz operators and the complex multiplications.
\end{abstract}

\noindent\textbf{Keywords:} Algebraic geometry (Abelian varieties, \'etale cohomology); 14F20, 14K22

\section*{1.}

Soit $k$ une cl\^oture alg\'ebrique du corps $\mathbf{F}_p$ \`a $p$ \'el\'ements. On d\'esignera par $\ell$ un nombre premier diff\'erent de $p$. On fixe, sur $k$, un isomorphisme $\mathbf{Q}_\ell \xrightarrow{\sim} \mathbf{Q}_\ell(1)$.

Soit $A$ une vari\'et\'e ab\'elienne sur $k$. On note $Z$, ou $Z(A)$, le groupe des cycles \`a coefficients rationnels sur $A$, gradu\'e par la codimension :
\[
Z = \bigoplus_i Z^i .
\]
On dispose sur $Z^i$ des relations d'\'equivalence num\'erique et d'\'equivalence cohomologique, pour la cohomologie \'etale \`a coefficients dans $\mathbf{Q}_\ell$. Deux cycles cohomologiquement \'equivalents sont num\'eriquement \'equivalents.

Soient $A_j$ les facteurs simples de $A$, $E_j$ le centre de $\End(A_j) \otimes \mathbf{Q}$ et $F$ le compos\'e, dans une cl\^oture alg\'ebrique de $\mathbf{Q}$, des cl\^otures galoisiennes des $E_j$; $F$ est donc un corps totalement r\'eel ou un corps CM. Si $\ell$ est un nombre premier, soit $D_\ell \subset \Gal(F/\mathbf{Q})$ un groupe de d\'ecomposition en $\ell$; soit $c \in \Gal(F/\mathbf{Q})$ la conjugaison complexe.

\begin{theorem}
Si $c \in D_\ell$, l'\'equivalence num\'erique et l'\'equivalence cohomologique, pour la cohomologie \'etale $\ell$--adique, co\"incident pour $A$.
\end{theorem}

Un th\'eor\`eme moins pr\'ecis, o\`u la condition sur $\ell$ d\'ependait, pour une vari\'et\'e simple, du choix d'un corps maximal de multiplication complexe, \'etait d\'emontr\'e dans [2].

Les remarques suivantes sont \'evidentes :

(a) $A$ \'etant donn\'ee, la condition du Th\'eor\`eme 1 est v\'erifi\'ee, d'apr\`es le th\'eor\`eme de \v{C}ebotarev, pour un ensemble $\mathcal{L}$ de nombres premiers $\ell$ de densit\'e (analytique ou na\"ive) non nulle.

(b) Si la condition du Th\'eor\`eme 1 est v\'erifi\'ee pour $A$, elle l'est aussi pour les puissances $A^n$ de $A$. Il existe donc un ensemble $\mathcal{L}$, de densit\'e $> 0$, tel que la conclusion soit vraie pour toute puissance $A^n$. Ceci avait \'et\'e remarqu\'e par Milne [5].

Soit $H = \bigoplus_{i=0}^{2g} H^i(A, \mathbf{Q}_\ell) = \bigoplus_{i=0}^{2g} H^i$, o\`u $g = \dim(A)$. Soit $D = \End(A) \otimes \mathbf{Q}$.

Le groupe $D^\times$ op\`ere sur $H$, ainsi que le groupe, isomorphe \`a $\SL(2, \mathbf{Q})$, donn\'e par la th\'eorie de Lefschetz. Tous deux pr\'eservent les classes de cohomologie alg\'ebriques. Soit $G \subset \GL(H)$ l'adh\'erence de Zariski du groupe engendr\'e par $D^\times$ et $\SL(2)$. Le point--cl\'e de notre d\'emonstration du Th\'eor\`eme 1 est une description de $G$, ou plut\^ot du groupe $G_{\mathbf{Q}_\ell}$ sur $\mathbf{Q}_\ell$ qui s'en d\'eduit par extension des scalaires. Nous montrerons au \S\,4 comment cette description se ram\`ene aux deux cas suivants.

Soient d'abord $E$ une courbe elliptique supersinguli\`ere et $A = E^n$. Le groupe $\End(A)^\times$ agissant sur $H \otimes \mathbf{Q}_\ell = \Lambda^\bullet(H^1 \otimes \mathbf{Q}_\ell)$ a pour adh\'erence de Zariski le groupe $\GL(H^1 \otimes \mathbf{Q}_\ell)$. Soit $V = H^1 \otimes \mathbf{Q}_\ell$ et $V^*$ son dual.

\begin{theorem}
Le groupe $G_{\mathbf{Q}_\ell}$ est isomorphe au groupe $\CSpin(V \oplus V^*)$ des similitudes spinorielles, $V \oplus V^*$ \'etant muni de la forme quadratique naturelle.

Il op\`ere sur $H \otimes \mathbf{Q}_\ell = \bigoplus_i H^i(A, \mathbf{Q}_\ell)$ par la somme de ses deux repr\'esentations demi--spinorielles. En particulier, la repr\'esentation sur les classes de degr\'e pair (resp.\ impair) est irr\'eductible.
\end{theorem}

\`A l'oppos\'e, soit $A$ une vari\'et\'e irr\'eductible telle que $\End(A) \otimes \mathbf{Q} = D$ soit une alg\`ebre \`a division, de degr\'e r\'eduit $n$, sur un corps quadratique imaginaire $E$; $A$ est de dimension $n$.

Soit $c \in \Gal(E/\mathbf{Q})$ la conjugaison complexe, et choisissons un plongement $j$ de $E$ dans $\mathbf{Q}_\ell$. La structure de $E$--module de $H^1(A, \mathbf{Q}_\ell)$ fournit une d\'ecomposition $H^1 = H^{10} \oplus H^{100}$, o\`u $\lambda \in E$ agit sur $H^{10}$ (resp.\ $H^{100}$) par multiplication par $j(\lambda)$ (resp.\ $j(c\lambda)$). Soit $V = V_1 \oplus V_2$ la d\'ecomposition duale de $H^1(A, \mathbf{Q}_\ell)$. L'adh\'erence de Zariski de $D^\times$ devient, sur $\mathbf{Q}_\ell$, le groupe $\GL(V_1) \times \GL(V_2)$.

Munissons
\[
H^\bullet(A, \mathbf{Q}_\ell) = \bigoplus_{p,q} \Lambda^p V_1 \otimes \Lambda^q V_2
\]
de la graduation donn\'ee par $k := p - q$ :
\begin{equation}\label{eq:grad}
\mathrm{gr}^k = \bigoplus_{p-q=k} \Lambda^p V_1 \otimes \Lambda^q V_2 .
\end{equation}
La partie de degr\'e $(-n)$ est donc $\Lambda^0 V_1 \otimes \Lambda^n V_2$, de dimension $1$, et celle de degr\'e $(-n+1)$ est $V_1 \otimes \Lambda^n V_2 \oplus \Lambda^0 V_1 \otimes \Lambda^{n-1} V_2$, de dimension $2n$.

\begin{theorem}\label{thm:3}
\emph{(i)} Pour un isomorphisme de $\bigoplus_k \Lambda^p V_1 \otimes \Lambda^q V_2$ avec $\Lambda^\bullet(\mathrm{gr}^{-n+1})$, transformant $\mathrm{gr}^k$ en $\Lambda^{\frac{n+k}{2}} \mathrm{gr}^{-n+1}$, $G_{\mathbf{Q}_\ell}$ s'identifie au groupe engendr\'e par $\GL(\mathrm{gr}^{-n+1})$ agissant sur $\Lambda^\bullet \mathrm{gr}^{-n+1}$, et par les homoth\'eties $x \mapsto \lambda x$.

\emph{(ii)} Les $\mathrm{gr}^k$ ($-n \leq k \leq n$) sont donc des repr\'esentations irr\'eductibles non isomorphes de $G$. Les sous--espaces $\mathrm{gr}^k$ et $\mathrm{gr}^\ell$ sont orthogonaux pour le cup--produit si $k + \ell \neq 0$, et en dualit\'e parfaite si $k + \ell = 0$.
\end{theorem}

\section*{2.}

Montrons comment r\'eduire le Th\'eor\`eme 1 (et m\^eme une variante un plus plus forte, cf.\ \S\,5) \`a l'analyse du groupe $G$.

(a) Soient $Z_0$ le $\mathbf{Q}$--sous--espace vectoriel de $H^\bullet(A, \mathbf{Q}_\ell)$ image de $Z$, et $Z_1$ le $\mathbf{Q}_\ell$--sous--espace vectoriel engendr\'e par $Z_0$. Nous montrerons que la forme d'intersection est non d\'eg\'en\'er\'ee sur $Z_1$. Sa restriction (qui est rationnelle) \`a $Z_0$ a donc un noyau nul : elle est non d\'eg\'en\'er\'ee et induit une surjection de $Z_1$ dans $\Hom(Z_0, \mathbf{Q}_\ell)$. On a donc $\dim_{\mathbf{Q}_\ell}(Z_1) \geq \dim_{\mathbf{Q}}(Z_0)$, les \'el\'ements d'une base de $Z_0$ sont lin\'eairement ind\'ependants sur $\mathbf{Q}_\ell$, et $Z_0 \otimes \mathbf{Q}_\ell \xrightarrow{\sim} Z_1$.

(b) On notera $\mathbf{Q}_\lambda$ une extension galoisienne finie assez grande de $\mathbf{Q}_\ell$. Il suffit alors de v\'erifier la non--d\'eg\'en\'erescence de la forme d'intersection sur $Z_\lambda := Z_1 \otimes_{\mathbf{Q}_\ell} \mathbf{Q}_\lambda \subset H_\lambda := H^\bullet(A, \mathbf{Q}_\lambda)$.

(c) Soit $L$ l'op\'erateur de Lefschetz, op\'erant sur $H$, d\'eduit d'une polarisation sur $A$, et soit $\Lambda : H^i \to H^{i-2}$ l'op\'erateur associ\'e [4]. Les exponentielles des endomorphismes nilpotents $L$ et $\Lambda$ de $H$ engendrent un sous--groupe $S$, isomorphe \`a $\SL(2, \mathbf{Q})$. Soit $G$ l'adh\'erence de Zariski, dans $\GL(H)$, du sous--groupe engendr\'e par $S$ et par $(\End(A) \otimes \mathbf{Q})^\times$. Il pr\'eserve $Z_1$, donc $Z_\lambda$ (le fait que $\Lambda$ pr\'eserve les cycles alg\'ebriques est d\^u \`a Lieberman; cf.\ [4]). Puisque $G$ contient le groupe $\mathbf{G}_m$ correspondant \`a la graduation de $H^\bullet(A)$ (action de $(\mathbf{Q} \cdot 1)^\times \subset \End(A) \otimes \mathbf{Q}$), ainsi que celui qui correspond \`a la m\^eme graduation, d\'ecal\'ee de $\dim(A)$ (matrices diagonales dans $S$), il contient les homoth\'eties $x \mapsto \lambda x$.

(d) Supposons que $\mathbf{Q}_\lambda$ est une extension galoisienne assez grande de $\mathbf{Q}_\ell$. On v\'erifiera que $H_\lambda$ est une somme de repr\'esentations absolument irr\'eductibles non isomorphes $W$ de $G$, et que la forme d'intersection d\'efinit des accouplements non d\'eg\'en\'er\'es sur $W \times W'$, $W \mapsto W'$ \'etant une involution des facteurs. Supposons qu'il existe $\sigma \in \Gal(\mathbf{Q}_\lambda/\mathbf{Q}_\ell)$ \'echangeant les facteurs en dualit\'e. Puisque $Z_\lambda$ est stable par le groupe de Galois, la forme d'intersection est alors non d\'eg\'en\'er\'ee : ceci terminera la d\'emonstration.

Soit $E$ le centre de $\End(A) \otimes \mathbf{Q}$ : $E$ est un produit de corps CM $E_i$ (on convient qu'un corps totalement r\'eel est CM\ldots). Supposons que $E$ a la propri\'et\'e suivante, \'equivalente \`a l'hypoth\`ese ``$c \in D_\ell$'' du Th\'eor\`eme 1 :
\begin{equation}\label{eq:hyp}
\text{Il existe une extension galoisienne $\mathbf{Q}_\lambda$ de $\mathbf{Q}_\ell$ et $\sigma \in \Gal(\mathbf{Q}_\lambda/\mathbf{Q}_\ell)$ tels}
\end{equation}
\[
\text{que, pour tout $i$ et pour tout plongement $\tau : E_i \to \overline{\mathbf{Q}}_\ell$, on ait $\sigma\tau = \tau \circ c$,}
\]
\[
c \text{ \'etant la conjugaison complexe sur } E_i .
\]

On v\'erifiera alors ci--dessous l'existence de $\sigma$ \'echangeant $W$ et $W'$. Sous cette hypoth\`ese, donc, le Th\'eor\`eme 1 s'ensuit.

\section*{3.}

Voici quelques rappels d'alg\`ebre lin\'eaire, pour lesquels nous renvoyons \`a [3], bas\'e sur [1]. Soit $V$ un espace vectoriel sur un corps $K$ (ici $\mathbf{Q}_\ell$, $\mathbf{Q}_\lambda$ ou $\mathbf{Q}_\ell$), $V^*$ son dual, et munissons $W = V \oplus V^*$ de la forme quadratique non d\'eg\'en\'er\'ee
\[
Q(v + \lambda) = \langle \lambda, v \rangle \qquad (v \in V, \lambda \in V^*) .
\]
Soit $C$ l'alg\`ebre de Clifford associ\'ee, quotient de l'alg\`ebre tensorielle $T(V \oplus V^*)$ par l'id\'eal engendr\'e par les $x \cdot x - Q(x)$ ($x \in V \oplus V^*$). C'est une alg\`ebre gradu\'ee par $\mathbf{Z}/2\mathbf{Z}$ [3]; soit $C = C^+ \oplus C^-$ cette graduation. Faisons agir $W$ sur l'alg\`ebre ext\'erieure $\Lambda^\bullet V$ de $V$ par
\begin{equation}\label{eq:clifford}
v \cdot x = v \wedge x \quad (v \in V, x \in \Lambda^\bullet V) \qquad
\lambda \cdot x = \lambda \rfloor x \quad (\lambda \in V^*, x \in \Lambda^\bullet V)
\end{equation}
o\`u $\lambda \rfloor x = i(\lambda)x$. La relation qui d\'efinit $C$ est v\'erifi\'ee, et \eqref{eq:clifford} se prolonge donc en une structure de $C$--module
\begin{equation}\label{eq:clifford-end}
C \longrightarrow \End(\Lambda^\bullet V) .
\end{equation}
Le morphisme \eqref{eq:clifford-end} est un isomorphisme : voir [1], Thm.~2; ou remarquer que si $V = V' \oplus V''$, $C$ est le produit tensoriel (au sens $\mathbf{Z}/2$--gradu\'e) des alg\`ebres analogues $C'$ et $C''$ : on est donc ramen\'e au cas, facile, o\`u $V$ est de rang 1.

L'image de $C^+$ par \eqref{eq:clifford-end} est l'alg\`ebre des endomorphismes de $\Lambda^\bullet V$ respectant la $\mathbf{Z}/2$--graduation de $\Lambda^\bullet V$.

Le groupe $C^\times$ agit sur l'alg\`ebre $C$ par conjugaison. Le sous--groupe $H$ qui normalise $W \subset C$ respecte la $\mathbf{Z}/2$--graduation de $C$, donc est form\'e d'\'el\'ements pairs ou impairs. Son action sur $W$ d\'efinit $H \to \Ogrp(W)$, de noyau le centre de $C^\times$. Le sous--groupe $(C^+)^\times \cap H$ de $(C^+)^\times$ qui normalise $W$ est le groupe des similitudes spinorielles $\CSpin(W)$. Il s'envoie sur $\SO(W)$ :
\[
1 \longrightarrow \mathbf{G}_m \longrightarrow \CSpin(W) \longrightarrow \SO(W) \longrightarrow 1 .
\]
Le groupe $\CSpin(W)$ est donc le groupe des automorphismes de $\Lambda^\bullet V$, respectant la $\mathbf{Z}/2$--graduation et le sous--espace $W$ de $\End(\Lambda^\bullet V)$. Son alg\`ebre de Lie est celle des endomorphismes pairs $T$ de $\Lambda^\bullet V$ tels que, pour $W$ identifi\'e \`a un sous--espace de $\End(\Lambda^\bullet V)$, on ait $[T, W] \subset W$. La repr\'esentation $\Lambda^\bullet V$ de $\CSpin(W)$ est la somme des deux repr\'esentations semi--spinorielles.

Il est clair que $\GL(V)$, agissant sur $\Lambda^\bullet V$, est un sous--groupe de $\CSpin(W)$. Il rel\`eve dans $\CSpin$ le sous--groupe de Levi de $\SO(W)$, intersection des sous--groupes paraboliques oppos\'es respectant les sous--espaces isotropes maximaux $V$ et $V^*$ de $W$.

Le sous--espace isotrope maximal $V^* \subset W \subset C$ et la droite $K \cdot 1$ de $\Lambda^\bullet V$ se d\'eterminent mutuellement : $V^*$ est le sous--espace de $W$ annulant $1$ et $K \cdot 1$ est le sous--espace de $\Lambda^\bullet V$ annul\'e par $V^*$. Par la m\^eme r\`egle, $V$ correspond \`a la droite $\Lambda^{\max} V$.

La m\^eme histoire se r\'ep\`ete pour tout sous--espace isotrope maximal de $W$. Le cas qui nous int\'eresse est celui o\`u $V = V_1 \oplus V_2$, et o\`u $W$ est d\'ecompos\'e en la somme de deux sous--espaces isotropes maximaux par
\[
W = (V_1 \oplus V_2^*) \oplus (V_1^* \oplus V_2) .
\]
L'annulateur de $V_1^* \oplus V_2$ est la droite $\Lambda^{\max} V_2 \subset \Lambda^\bullet V$. Fixons--en un g\'en\'erateur $\Omega$. L'inclusion dans $W$ du sous--espace isotrope compl\'ementaire $V_1 \oplus V_2^*$ se prolonge en un morphisme d'alg\`ebres
\[
\Lambda^\bullet(V_1 \oplus V_2^*) \longrightarrow C ,
\]
et
\begin{equation}\label{eq:isom}
I : \Lambda^\bullet(V_1 \oplus V_2^*) \longrightarrow \Lambda^\bullet V, \qquad \lambda \longmapsto \lambda \cdot \Omega
\end{equation}
est un isomorphisme. Le super--crochet avec $x \in W$ stabilise $\Lambda^\bullet(V_1 \oplus V_2^*) \subset C$ et induit une super--d\'erivation. Pour $x \in V_1^* \oplus V_2$, identifi\'e au dual de $V_1 \oplus V_2^*$, cette d\'erivation est le produit int\'erieur $x \rfloor$. Via l'isomorphisme $I$, l'action de $v \in V_1 \oplus V_2^* \subset W$ devient donc un produit ext\'erieur, et celle de $x \in V_1^* \oplus V_2 \subset W$, qui annule $\Omega$, devient un produit int\'erieur.

Le produit ext\'erieur dans $\Lambda^\bullet(V_1 \oplus V_2^*)$, transport\'e par $I$, d\'efinit un nouveau produit, not\'e $*$, sur $\Lambda^\bullet V$. Si on munit $\Lambda^\bullet(V_1 \oplus V_2^*)$ de la graduation pour laquelle $V_1$ (resp.\ $V_2^*$) est de degr\'e $1$ (resp.\ $-1$), le produit est compatible avec la graduation. L'isomorphisme $I$ la transporte en la graduation \'evidente de $\Lambda^\bullet V$, translat\'ee par $\dim(V_2)$ pour placer $\Omega$ en degr\'e $0$. Le produit $*$ est donc compatible avec la graduation \'evidente translat\'ee. La graduation naturelle de $\Lambda^\bullet(V_1 \oplus V_2^*)$ (pour laquelle $V_2^*$ est de degr\'e $1$) est aussi compatible au produit. On obtient alors sur $\Lambda^\bullet V$ la graduation $\mathrm{gr}^k$ de \eqref{eq:grad}, d\'ecal\'ee de $\dim(V_2)$ ($1 \in \Lambda^\bullet V_2$, de degr\'e $0$ pour $\mathrm{gr}^k$, est de degr\'e $\dim(V_2)$).

Pour $v_1 \in V_1$, $v_2 \in V_2$, calculons $(v_1 \wedge v_2)^\wedge : \Lambda^\bullet V \to \Lambda^\bullet V$ \`a l'aide de l'isomorphisme \eqref{eq:isom} : on a le diagramme
\[
\begin{array}{ccc}
\Lambda^\bullet(V_1 \oplus V_2^*) & \longrightarrow & \Lambda^\bullet(V_1 \oplus V_2^*) \\
\downarrow\rlap{$\scriptstyle I$} & & \downarrow\rlap{$\scriptstyle I$} \\
\Lambda^\bullet V & \longrightarrow & \Lambda^\bullet V
\end{array}
\]
et l'application du bas est un produit tensoriel d'applications. Si $a \in \Lambda^\bullet V_1$, $b \in \Lambda^\bullet V_2$,
\[
(v_1 \wedge v_2)(a \wedge b) = (-1)^{\deg a} (v_1 \wedge a) \wedge (v_2 \wedge b) .
\]
Pour l'isomorphisme $\Lambda^\bullet V_2^* \to \Lambda^\bullet V_2$ d\'efinissant $I$, $b \mapsto v_2 \wedge b$ se transporte en $b^* \mapsto v_2 \rfloor b^*$. On en d\'eduit que $(v_1 \wedge v_2)^\wedge$ se transporte en l'endomorphisme $v_1 \wedge v_2 \rfloor$ de $\Lambda^\bullet(V_1 \oplus V_2^*)$ dont on v\'erifie aussit\^ot que c'est une d\'erivation. Celle--ci respecte $V_1 \oplus V_2^*$.

\section*{4.}

Calculer le groupe $G$ (du \S\,1, et sur une extension assez grande de $\mathbf{Q}_\ell$) est amusant. Voici l'histoire.

(a) Soit $E$ le centre de $D = \End(A) \otimes \mathbf{Q}$, produit de corps CM $E_i$; $D$ est centrale simple sur $E$.

(b) Apr\`es extension des scalaires de $\mathbf{Q}$ \`a $\mathbf{Q}_\lambda$ (assez grand), $E$ devient une somme de copies de $\mathbf{Q}_\lambda$, index\'ees par l'ensemble $\Sigma$ des morphismes de $\mathbf{Q}$--alg\`ebres $E \to \mathbf{Q}_\lambda$. Apr\`es extension des scalaires de $\mathbf{Q}_\ell$ \`a $\mathbf{Q}_\lambda$, $H^1(A)$ devient $\bigoplus H_\sigma$ ($\sigma \in \Sigma$), o\`u
\[
H_\sigma = \{\alpha \in H^1(A, \mathbf{Q}_\lambda) : [z]\alpha = \sigma(z)\alpha \ (z \in E)\} ,
\]
$[z]$ d\'esignant l'action de $E$ sur $H^1(A)$ d\'eduite de son action sur $A$. L'alg\`ebre $D \otimes \mathbf{Q}_\lambda$ est $\prod_\sigma \End(H_\sigma)$.

Le groupe $G$, qui agit sur $\Lambda^\bullet\bigl(\bigoplus_\sigma H_\sigma\bigr)$, est engendr\'e par le produit des $\GL(H_\sigma)$ et par un $\SL(2)$ ``diagonal''. Soit $L : H^i(A, \mathbf{Q}_\lambda) \to H^{i+2}(A, \mathbf{Q}_\lambda)$ l'op\'erateur de Lefschetz, que l'on identifie \`a une classe de cohomologie dans $H^2(A, \mathbf{Q}_\lambda) = \Lambda^2\bigl(\bigoplus H_\sigma\bigr)$; le Lemme 2.1 de [2] implique que
\[
L = \sum_v L_v
\]
o\`u $v$ parcourt l'ensemble quotient $\Sigma/\{1, c\}$, la conjugaison complexe op\'erant sur $\Sigma$ par son action sur $E$, et o\`u
\[
L_v \in \Lambda^2(H_\sigma \oplus H_{c\sigma})
\]
pour $v = \{\sigma, c\sigma\}$. ($\sigma$ peut \^etre fix\'e par $c$, auquel cas $L_v \in \Lambda^2 H_\sigma$). Chaque $L_v$ est dans l'adh\'erence de l'orbite de $L$ sous $\prod \GL(H_\sigma)$ et donc dans l'alg\`ebre de Lie de $G$.

Les m\^emes arguments donnent une d\'ecomposition analogue de l'op\'erateur $\Lambda$; les relations de commutation sont pr\'eserv\'ees par les composantes $L_v$ et $\Lambda_v$. Par cons\'equent, $G$ est aussi engendr\'e par le produit des $\GL(H_\sigma)$ et d'une copie de $\SL(2)$ d\'efinie sur $\mathbf{Q}_\lambda$ pour chaque orbite de $c$ dans $\Sigma$. Ainsi $G$ est un produit agissant sur un produit tensoriel.

(c) \textit{Histoire pour $\{\sigma\}$, $\sigma = c\sigma$.}

Ceci correspond \`a un \'eventuel facteur $\mathbf{Q}$ de $E$, donc \`a un facteur isog\`ene \`a $B = E^n$ de $A$, o\`u $E$ est une courbe elliptique supersinguli\`ere sur $k$.

Soit $V = H_\sigma = H^1(B, \mathbf{Q}_\lambda)$, espace de dimension paire. Soit $W = V \oplus V^*$. La classe de Lefschetz $L \in \Lambda^2 V$ est un isomorphisme de $V$ vers $V^*$. On va montrer que le groupe engendr\'e par $\GL(V)$ et le sous--groupe $\SL(2)$ associ\'e \`a $L$ est $\CSpin(W)$ dans sa repr\'esentation spinorielle sur $H^\bullet(B) = \Lambda^\bullet V$.

On v\'erifie aussit\^ot que, pour son action sur $\Lambda^2 V$ par multiplication \`a gauche, $L$ v\'erifie
\[
[L, W] \subset W .
\]
Par cons\'equent, $L$ est un \'el\'ement de $\mathrm{Lie}(\CSpin(W))$; celui--ci appartient \`a l'alg\`ebre de Lie du radical unipotent $N$ du sous--groupe de Levi d\'efini par $W = V \oplus V^*$ (\S\,3) et en est un \'el\'ement non--singulier. Les conjugu\'es par $\GL(V)$ de l'exponentielle de $L$ engendrent donc $N$. De m\^eme, $\Lambda$ et ses conjugu\'es engendrent l'alg\`ebre de Lie du sous--groupe unipotent oppos\'e. Par cons\'equent, le groupe engendr\'e par $\GL(V)$ et $\SL(2)$ contient le groupe de $\Spin$; on a d\'ej\`a vu qu'il contenait les homoth\'eties (\S\,2, (c)). Son adh\'erence de Zariski, sur $\mathbf{Q}_\ell$, est donc $\CSpin(V \oplus V^*)$, ce qui d\'emontre le Th\'eor\`eme 2.

(d) \textit{Histoire pour $\{\sigma, c\sigma\}$, $\sigma \neq c\sigma$.}

Dans ce cas $H_v = H_\sigma \oplus H_{c\sigma} = V_1 \oplus V_2$ et $(\End(A) \otimes \mathbf{Q}_\lambda)^\times$ op\`ere sur $H_v$ par $\GL(V_1) \times \GL(V_2)$.

Ecrivons, selon le \S\,3 :
\[
\Lambda^\bullet(V_1 \oplus V_2) \xleftarrow{\ I^{-1}\ } \Lambda^\bullet(V_1 \oplus V_2^*) .
\]
On munit $\Lambda^\bullet(V_1 \oplus V_2)$ du produit $*$, qui devient par $I^{-1}$ le produit ext\'erieur \`a droite. On peut \'ecrire ([2, Lemme 2.1]) :
\[
L = \sum e_i \wedge f_i \in \Lambda^2(V_1 \oplus V_2)
\]
avec $e_i \in V_1$, $f_i \in V_2$. L'isomorphisme $I^{-1}$ \'echange l'action de $L$ sur $\Lambda^\bullet(V_1 \oplus V_2)$ et l'action sur le membre de droite donn\'ee par
\[
\sum e_i \wedge f_i \rfloor .
\]
Celle--ci est une d\'erivation (\S\,3); en fait, dans les bases de $V_1$ et $V_2^*$ donn\'ees par $(e_i, f_i^*)$, c'est la somme des endomorphismes
\[
x = \begin{pmatrix} 0 & 1 \\ 0 & 0 \end{pmatrix} \in \mathfrak{sl}(2)
\]
de $k^2$, donc l'endomorphisme de $k^{2n}$ correspondant \'etendu (comme action d'alg\`ebre de Lie) \`a $\Lambda^\bullet(V_1 \oplus V_2^*) = \Lambda^\bullet(k^{2n})$. L'action de $L$ sur $\Lambda^\bullet(V_1 \oplus V_2)$ laisse donc invariant (au sens des alg\`ebres de Lie) le produit $*$, c'est donc une d\'erivation.

Par ailleurs l'\'el\'ement $h = \begin{pmatrix} 1 & 0 \\ 0 & -1 \end{pmatrix}$ de $\mathfrak{sl}(2)$ op\`ere sur $\Lambda^i V$ par $i - n = \mathrm{gr}^1(\Lambda^i V)$ o\`u $\mathrm{gr}^1$ est la graduation naturelle translat\'ee (\S\,3). Celle--ci \'etant compatible au produit $*$, celui--ci est invariant par la sous--alg\`ebre de Borel de $\mathfrak{sl}(2)$. Or dans une repr\'esentation de $\mathfrak{sl}(2)$, un vecteur annul\'e par une sous--alg\`ebre de Borel l'est par $\mathfrak{sl}(2)$. On en d\'eduit que $*$ est invariant par $\SL(2)$.

Consid\'erons la d\'ecomposition de $\Lambda^\bullet V$ donn\'ee par \eqref{eq:grad} :
\[
\mathrm{gr}^k \Lambda^\bullet V = \bigoplus_{p-q=k} \Lambda^p V_1 \otimes \Lambda^q V_2 .
\]
Noter que cette d\'ecomposition est invariante par $L$ et $\Lambda$. En effet, $E$ (par son plongement dans $\mathbf{Q}_\lambda$) op\`ere sur $\Lambda^\bullet V$. Soit $\mathrm{U}(1)$ le groupe $\{\lambda \in E : N_{E/F} \lambda = 1\}$ o\`u $F$ est la sous-alg\`ebre de $E$ form\'ee des points fixes de $c$. Si $\lambda \in \mathrm{U}(1)$, $\lambda$ op\`ere sur $\mathrm{gr}^k$ par le scalaire $\lambda^{p-q} (\in \mathbf{Q}_\lambda)$; il comute \`a $L$ et $\Lambda$, et les $\lambda^{p-q}$ sont des caract\`eres distincts de $\mathrm{U}(1)$. R\'esumons :

\begin{lemma}\label{lem:4.1}
On munit $\Lambda^\bullet V$ du produit $*$, et de la graduation \eqref{eq:grad}.

\emph{(i)} $\Lambda^\bullet V$ est engendr\'e par $\mathrm{gr}^{1-n} = V_1 \otimes \Lambda^n V_2 \oplus \Lambda^0 V_1 \otimes \Lambda^{n-1} V_2$.

\emph{(ii)} Le produit $* : \mathrm{gr}^i \times \mathrm{gr}^j \to \mathrm{gr}^{i+j}$ est invariant par $\SL(2)$, $\SL(V_1)$ et $\SL(V_2)$.

\emph{(iii)} La graduation de $\Lambda^\bullet V$ est pr\'eserv\'ee par $G$.
\end{lemma}

L'image de $G$ dans $\GL(\mathrm{gr}^{n-1})$ contient $\GL(V_1)$, $\GL(V_2) \simeq \GL(\Lambda^{n-1} V_2)$ ainsi qu'un isomorphisme $\Lambda^0 V_1 \otimes \Lambda^n V_2 \xrightarrow{\sim} V_1 \otimes \Lambda^n V_2$ (donn\'e par $L$) et un isomorphisme inverse (donn\'e par $\Lambda$). C'est donc $\GL(\mathrm{gr}^{n-1})$. La propri\'et\'e (i) identifie $\Lambda^\bullet V$ \`a $\Lambda^\bullet(\mathrm{gr}^{n-1})$, de fa\c{c}on compatible \`a l'action de $\SL(2) \times \SL(V_1) \times \SL(V_2)$. Le groupe engendr\'e par $\{zs_1, z^{-1}s_2\}$, o\`u $z$ est un scalaire et $s_i \in \SL(V_i)$, dans $\GL(V_1) \times \GL(V_2)$ op\`ere sur $\Lambda^\bullet V$ par l'action de $\GL(\mathrm{gr}^{1-n})$ d\'ecrite par le Th\'eor\`eme 3. Le groupe $G$ est donc engendr\'e par $\GL(\mathrm{gr}^{1-n})$, isomorphe \`a $\GL(2n)$, et par un tore de dimension $1$, qui op\`ere par homoth\'eties globales sur $\Lambda^\bullet V$, d'o\`u le Th\'eor\`eme 3 (i). Il est clair que, d\'ej\`a pour l'action de $\GL(\mathrm{gr}^{1-n})$, les repr\'esentations obtenues sont absolument irr\'eductibles et non isomorphes; enfin, si $k + \ell \neq 0$, les espaces $\mathrm{gr}^k$ et $\mathrm{gr}^\ell$ ne sont pas des repr\'esentations duales, m\^eme pour l'action de $\GL(2n) \subset G$. Le cup--produit \'etant non d\'eg\'en\'er\'e, $\mathrm{gr}^k$ et $\mathrm{gr}^\ell$ sont donc en dualit\'e parfaite.

\section*{5.}

Une fois acquise la description de $G$, voici la d\'emonstration du Th\'eor\`eme 1.

Ecrivons
\[
H_\lambda = H^\bullet(A, \mathbf{Q}_\lambda) = \bigotimes_{v \in \Sigma/c} H_v^\bullet
\]
o\`u $H_v^\bullet$ est comme dans le \S\,4. On suppose $\mathbf{Q}_\lambda$ assez grand, et galoisien sur $\mathbf{Q}_\ell$.

Sous l'action du produit (d\'esormais not\'e $G$) des groupes d\'ecrits dans le \S\,4, $H_\lambda$ se d\'ecompose comme somme directe index\'ee par les fonctions $\sigma \mapsto k(\sigma)$ ($\sigma \in \Sigma$) o\`u
\begin{align}
k(\sigma) &\in \mathbf{Z} \text{ si } \sigma \neq \sigma c \text{; alors } k(\sigma c) = -k(\sigma) \label{eq:k1}\\
k(\sigma) &\in \mathbf{Z}/2\mathbf{Z} \text{ si } \sigma = \sigma c \text{ (facteur supersingulier \'eventuel)} . \label{eq:k2}
\end{align}

Dans \eqref{eq:k1}, $k(\sigma)$ est donn\'e par le Th\'eor\`eme \ref{thm:3}, \'etendu, comme dans le \S\,4, au cas o\`u $E$ est CM plut\^ot que quadratique imaginaire. On remarquera que si on change $\sigma$ en $\sigma c$, les espaces $V_1$ et $V_2$ du \S\,4 sont \'echang\'es, d'o\`u la relation \eqref{eq:k1}.

Soit $K$ le groupe des fonctions v\'erifiant \eqref{eq:k1} et \eqref{eq:k2}. Puisque l'espace vectoriel sur $\mathbf{Q}_\lambda$ engendr\'e par les cycles alg\'ebriques est stable par $G$ ainsi que par cup--produit, il correspond \`a un sous--mono\"ide $\mathrm{Alg}$ de $K$. De plus,
\begin{align}
&\mathrm{Alg} \subset K_0 = \{k : \sum_{\sigma \mod c} k(\sigma) \text{ pair}\} \label{eq:alg1}\\
&H_\lambda(k) \text{ est orthogonal \`a } H_\lambda(\ell) \text{ si } k + \ell \neq 0 \text{;} \label{eq:alg2}\\
&H_\lambda(k) \text{ et } H_\lambda(-k) \text{ sont en dualit\'e} \label{eq:alg3}\\
&\mathrm{Alg} \text{ est stable par } \Gal(\mathbf{Q}_\lambda/\mathbf{Q}_\ell) . \label{eq:alg4}
\end{align}

Supposons que le groupe de Galois contienne un \'el\'ement agissant comme $(-1)$ sur $K$. Alors, sur tout sous--espace $\mathbf{Q}_\ell$--rationnel de $H_\lambda$, invariant par $G$, la dualit\'e donn\'ee par la forme d'intersection est non d\'eg\'en\'er\'ee. C'est en particulier le cas si (avec les notations du \S\,1) la conjugaison complexe $c \in \Gal(E/\mathbf{Q})$ appartient au groupe de d\'ecomposition de $\ell$, d'o\`u le Th\'eor\`eme 1.

Un peu moins suffit : pour que tout sous--mono\"ide stable soit un sous--groupe de $K$ (donc que l'espace des cycles alg\'ebriques contienne, avec l'espace associ\'e \`a $k$, son dual associ\'e \`a $-k$), il suffit que $K \otimes \mathbf{Q}$ n'admette pas de droite fixe sous l'action de $\Gal(\mathbf{Q}_\lambda/\mathbf{Q}_\ell)$. En d'autres termes, si $\ell$ est non ramifi\'e dans les facteurs de $E$, il suffit que $\mathrm{Frob}_\ell$ agissant sur $\Sigma$ v\'erifie :
\[
\text{Pour tout } \sigma \in \Sigma\text{, il existe } a \in \mathbf{Z} : \mathrm{Frob}_\ell^a(\sigma) = \sigma c .
\]

\section*{R\'ef\'erences}

\begin{enumerate}[label={[\arabic*]}]
\item N.\ Bourbaki, Alg\`ebres de Clifford, in \textit{Alg\`ebre}, livre II, Ch.~IX, \S\,9, Hermann, Paris 1959.
\item L.\ Clozel, Equivalence num\'erique et \'equivalence cohomologique pour les vari\'et\'es ab\'eliennes sur les corps finis, \textit{Ann.\ Math.}\ 150 (1999), 151--163.
\item P.\ Deligne, Clifford modules, in \textit{Quantum Fields and Strings, A course for Mathematicians}, Deligne et al.\ eds., vol.~I, AMS/IAS 1999, 107--112.
\item S.L.\ Kleiman, Algebraic cycles and the Weil conjectures, in \textit{Dix expos\'es sur la cohomologie des sch\'emas}, North--Holland, Amsterdam, Paris, 1968, 359--386.
\item J.S.\ Milne, The Tate conjecture for certain Abelian varieties over finite fields, \textit{Acta Arith.}\ 100 (2001), 135--166.
\end{enumerate}

\vspace{1em}

\noindent
Pierre Deligne\\
School of Mathematics, Institute for Advanced Study\\
Princeton NJ 08540 (USA)\\
E-mail : \texttt{deligne@math.ias.edu}

\vspace{0.5em}

\noindent
Laurent Clozel\\
Math\'ematiques, B\^at.~425, Universit\'e Paris--Sud, 91405 Orsay Cedex, France\\
E-mail : \texttt{Laurent.Clozel@math.u-psud.fr}

\end{document}
